% !TEX program = lualatex
\documentclass[a4paper,11pt]{scrartcl}

\usepackage[
  lang=zh,
  mode=journal,
  theme=ocean,
  toc=true,
  toc-layout=page
]{synaptic}
\usepackage{booktabs,tabularx}
\newcolumntype{Y}{>{\raggedright\arraybackslash}X}
\newcommand{\DocCmd}[1]{\textcolor{syn-primary!80!black}{\texttt{\textbackslash#1}}}

\SynapticNewTheorem[remark]{synobservation}{观察}
\SynapticHeader{synaptic 用户手册}
\SynapticSubHeader{面向 LuaLaTeX 的稳定学术排版}
\SynapticPreprint{v5.0.0}
\SynapticTitle{synaptic 宏包用户手册}
\SynapticAuthor{synaptic 项目}
\SynapticAffiliation{开源 LaTeX 社区}
\SynapticEmail{synaptic@example.com}
\SynapticSubject{synaptic 5.0.0 用户手册}
\SynapticKeywords{synaptic, LaTeX3, LuaLaTeX, 学术排版}

\begin{synabstract}
  \textbf{synaptic} 是一套模块化 LaTeX3 学术文档设计系统，适用于期刊
  论文、书籍、讲义与笔记。本手册介绍它的安装方式、环境要求、安全配置
  模型、排版系统、元数据 API、定理体系、主题与自定义工具以及各模式功能。
\end{synabstract}

\begin{document}
\SynapticMakeTitle

\section{快速开始}

宏包要求 LuaLaTeX 和 KOMA-Script 文档类。宏包已安装到 TeX 目录树时，
文档只需：

\begin{verbatim}
\documentclass{scrartcl}
\usepackage[mode=journal]{synaptic}
\SynapticTitle{一篇短文}
\SynapticAuthor{作者}
\begin{synabstract}摘要。\end{synabstract}
\begin{document}
\SynapticMakeTitle
\section{引言}
正文。
\end{document}
\end{verbatim}

书籍模式应使用 \texttt{scrbook}。纸张尺寸由文档类决定，synaptic 不会强制 A4。

\subsection{不安装直接使用仓库副本}

先提取宏包模块，再将 \texttt{TEXINPUTS} 指向构建输出：

\begin{verbatim}
l3build unpack
export TEXINPUTS=/path/to/synaptic/build/unpacked//:$TEXINPUTS
lualatex document.tex
\end{verbatim}

\subsection{安装到系统}

TeX Live 用户可在仓库根目录安装（版本更新时可直接替换）：

\begin{verbatim}
l3build install     # 复制到用户 texmf 目录树
l3build uninstall   # 移除旧版本
l3build ctan        # 生成 TDS/CTAN 压缩包
\end{verbatim}

任何标准 \texttt{texmf} 目录树均可；生成文件应位于
\texttt{tex/latex/synaptic/}。\texttt{l3build ctan} 生成的压缩包可直接解压到
同样的目录结构。

\subsection{环境要求}

\begin{itemize}
  \item LuaLaTeX；pdfLaTeX 与 XeLaTeX 会被明确拒绝。
  \item KOMA-Script 文档类（\texttt{scrartcl}、\texttt{scrreprt} 或
    \texttt{scrbook}）。
  \item 较新的 TeX 发行版，包含 \texttt{expl3}、\texttt{fontspec}、
    \texttt{unicode-math}、\texttt{geometry}、\texttt{microtype}、
    \texttt{thmtools} 等标准依赖。
  \item 系统中安装 Libertinus Serif、Sans、Mono 与 Math 字体。
  \item 教学与笔记卡片需要 \texttt{tcolorbox} 与 PGF；\texttt{lang=zh}
    还需要 \texttt{ctex}、思源宋体/黑体 CN 与霞鹜文楷。
\end{itemize}

\section{配置}

\noindent{\setlength{\parindent}{0pt}\noindent\begin{tabularx}{\textwidth}{@{}l l Y@{}}
  \toprule
  \textbf{键} & \textbf{默认值} & \textbf{可选值与作用} \\
  \midrule
  \texttt{mode} & journal & journal、book、lecture 或 notes \\
  \texttt{theme} & ocean & ocean、graphite、forest、midnight 或 paper \\
  \texttt{lang} & en & en 或 zh，决定标签语言与中文支持 \\
  \texttt{toc} & false & 标题流程是否输出目录 \\
  \texttt{toc-layout} & top & top、inline 或 page \\
  \texttt{title-layout} & auto & inline、compact、page、sequence 或模式默认值 \\
  \texttt{numbering} & auto & section、chapter、continuous、none 或模式默认值 \\
  \texttt{color-mode} & screen & screen、print 或 mono \\
  \texttt{density} & balanced & compact、balanced 或 airy \\
  \texttt{measure} & auto & narrow、standard、wide 或模式默认值 \\
  \texttt{theorem-style} & boxed & boxed 或 plain \\
  \texttt{binding-offset} & 0pt & 书籍装订补偿尺寸 \\
  \bottomrule
\end{tabularx}}

仅 \texttt{theme}、\texttt{toc} 和 \texttt{toc-layout} 可在加载后安全修改：

\begin{verbatim}
\SynapticSetup{theme=paper,toc=false,toc-layout=page}
\end{verbatim}

模式、语言与其他加载期排版选项会被冻结。运行时通过
\texttt{\string\SynapticSetup} 修改这些键会直接报错，而不是留下只完成一半的配置。

\subsection{模式}

\noindent{\setlength{\parindent}{0pt}\noindent\begin{tabularx}{\textwidth}{@{}l l Y@{}}
  \toprule
  \textbf{模式} & \textbf{建议文档类} & \textbf{用途} \\
  \midrule
  journal & scrartcl & 内联论文标题、元数据、运行页眉与声明 \\
  book & scrbook & 书籍扉页、篇章、镜像导航与前后文 \\
  lecture & scrartcl/scrreprt & 课程信息条与语义教学系统 \\
  notes & scrartcl & 紧凑信息条与轻量知识卡 \\
  \bottomrule
\end{tabularx}}

\textbf{journal} 在第一页渲染内联论文标题。运行页眉显示第一作者与短标题，
标题区支持摘要、关键词、arXiv、通讯、出版信息与投稿声明。

\textbf{book} 依次渲染半扉页、正式扉页、版权页与可选献词页，并提供前后文、
部分页、章副标题、题词与镜像双面页眉辅助命令。

\textbf{lecture} 渲染课程信息条（课程代码、标题、教师、学期、讲次），并以
简洁卡片系统包裹定理与例题。学习目标与课程小结界定每个教学单元。

\textbf{notes} 渲染紧凑信息条（可选标签行），支持文档类级双栏排版与末页
栏平衡。

\subsection{字体}

Synaptic 在所有机器上固定使用同一套字体：西文正文与数学使用 Libertinus Serif、
Sans、Mono 和 Math；Libertinus Math 没有独立粗体，粗体数学由同一字面确定性合成；
简体中文使用思源宋体/黑体 CN；中文强调使用霞鹜文楷。
缺少字体会直接报配置错误，绝不静默切换设计。

\section{排版系统}

synaptic 的默认外观是确定的而非装饰性的。以下设计原则贯穿所有模式。

\subsection{色彩交付}

{\setlength{\parindent}{0pt}\noindent\begin{tabularx}{\textwidth}{@{}l l Y@{}}
  \toprule
  \textbf{color-mode} & \textbf{用途} & \textbf{效果} \\
  \midrule
  screen & 屏幕阅读 & 完整主题色，彩色链接 \\
  print & 纸面输出 & 主题色，暗色（适合打印）链接 \\
  mono & 单色交付 & 仅黑灰令牌；语义色映射为灰度 \\
  \bottomrule
\end{tabularx}}

\subsection{密度与栏宽}

\textbf{density} 分三档调节行距：\texttt{compact} 适合密集资料汇编，
\texttt{balanced} 为按模式调整的默认值，\texttt{airy} 适合讲义与幻灯片式
材料。\textbf{measure} 覆盖模式默认栏宽：\texttt{narrow} 适合长篇正文，
\texttt{standard} 适合期刊排版，\texttt{wide} 适合设计驱动版面。书籍模式下
\texttt{measure} 调整镜像内外边距，保持双面版式对称。

\subsection{标题构成}

\textbf{title-layout} 决定各模式标题区的渲染方式。\texttt{auto} 解析为
journal 的 \texttt{inline}、book 的 \texttt{sequence} 以及 lecture 与 notes 的
\texttt{compact}。\texttt{page} 生成整页标题；\texttt{inline} 强制第一页
内联构成。所有标题渲染器共用一组非对称“连接轨道”：短主色线、强调节点与
淡色延伸线。书籍将其用于标题页和版权页的建筑性边缘，讲义与笔记则将其与
紧凑信息带组合，期刊模式保持内联且具有编辑感。

\subsection{导航}

所有模式在运行页眉下方绘制一条发丝线。发丝线使用独立的装饰性浅色令牌，
使导航保持安静。运行页眉与外侧页码重复微型节点标记；书籍模式还为双面强制空白页
提供克制的页脚锚点。

\subsection{标题、目录与图表标题}

所有编号的章节标题共用无衬线家族：\texttt{section} 使用主色，
\texttt{subsection} 使用深色正文色，\texttt{subsubsection} 更小且略微淡化，
使结构一目了然。启用目录时，顶层条目以加粗主色重现将页面标题颜色，深层条目
保持深色，并收紧编号对齐与垂直节奏。图表标题标签使用次要主题色，与整体设计
保持一致。

\subsection{卡片与规则}

教学与笔记卡片共享同一视觉语言：极小圆角、着语义色（主色、次级、警告或成功色）
的左侧色条、淡发丝线、浅色表面与可读标题颜色。规则与发丝线由主题主色的低浓度混合
推导，装饰永不与正文争抢。

\section{标题与 PDF 元数据}

四种模式共享元数据，但采用不同标题渲染器：

\noindent{\setlength{\parindent}{0pt}\noindent\begin{tabularx}{\textwidth}{@{}l Y@{}}
  \toprule
  \textbf{命令} & \textbf{含义} \\
  \midrule
  \DocCmd{SynapticTitle} & 必填标题与 PDF 标题 \\
  \DocCmd{SynapticSubtitle} & 可选副标题 \\
  \DocCmd{SynapticShortTitle} & 运行页眉短标题 \\
  \DocCmd{SynapticDate} & 显式日期 \\
  \DocCmd{SynapticAuthor} & 作者及可选机构标记，可重复 \\
  \DocCmd{SynapticAffiliation} & 机构，可重复 \\
  \DocCmd{SynapticEmail} & 可点击联系方式，可重复 \\
  \DocCmd{SynapticSubject} & PDF 主题 \\
  \DocCmd{SynapticKeywords} & 页面关键词与 PDF 关键词 \\
  \DocCmd{SynapticHeader} & 投稿抬头 \\
  \DocCmd{SynapticSubHeader} & 次级抬头 \\
  \DocCmd{SynapticPreprint} & 预印本编号 \\
  \DocCmd{SynapticArxiv} & arXiv 链接 \\
  \DocCmd{SynapticDedicated} & 献词 \\
  \DocCmd{SynapticMakeTitle} & 验证并输出标题信息 \\
  \bottomrule
\end{tabularx}}

journal 支持出版信息、通讯信息与作者注；book 支持半扉页、出版社、版次、ISBN 与
版权；lecture 支持课程、课程代码、讲次、教师与学期。论文的基金、数据、伦理或
利益冲突声明可用 \verb|\SynapticStatement{标题}{正文}|，其星号形式不写入目录。
PDF 字符串只进行一次 Unicode 安全转换。

\section{陈述与数学}

内置 \texttt{syntheorem}、\texttt{synlemma}、\texttt{synproposition}、
\texttt{syndefinition}、\texttt{syncorollary}、\texttt{synaxiom}、\texttt{synconjecture}、\linebreak
\texttt{synclaim}、\texttt{synproperty}、\texttt{syncriterion}、\texttt{synconvention}、\linebreak
\texttt{synproblem}、\texttt{synsolution}、\texttt{synexample}、\texttt{synexercise}、
\texttt{synremark}、\texttt{synnote} 与 \texttt{synwarning} 共享一个计数器。书籍默认按
章编号，短文档默认按节编号。定义正文使用正体；中文命题类正文保持正体，避免
伪斜体造成的阅读干扰。例题与练习为 plain 样式陈述，警告使用注记呈现。

默认（\texttt{theorem-style=boxed}）下，编号的陈述环境以精致的卡片呈现：微
妙的主题色背景、语义西侧强调条、淡发丝线与克制圆角，与教学与笔记卡片语言保持一致。
\texttt{synproof} 保持干净、不加框，同时不修改标准 \texttt{proof} 环境。
选择 \texttt{theorem-style=plain} 可恢复无框经典样式。

\begin{syntheorem}[高斯积分]
  \[
    \int_0^\infty e^{-x^2}\,\mathrm{d}x=\frac{\sqrt{\pi}}{2}.
  \]
\end{syntheorem}

\begin{synproof}
  对积分求平方，转换为极坐标，再计算径向积分。
\end{synproof}

自定义陈述可选择 \texttt{plain}、\texttt{definition} 或 \texttt{remark} 样式，
并会自动加框。环境名必须是以 \texttt{syn} 开头的小写字母数字标识符；
未命名空间化或重名都会报错。

\begin{verbatim}
\SynapticNewTheorem[remark]{synobservation}{观察}
\end{verbatim}

\begin{synobservation}
  可直接使用 unicode-math 提供的 \(\symbb{R}\)、\(\symcal{F}\) 与
  \(\symfrak{g}\)，无需额外生成单字母别名。
\end{synobservation}

\section{主题与自定义}

通过 \verb|\SynapticSetup{theme=名称}| 切换主题。先使用 xcolor 颜色表达式定义
不可变的自定义主题，再显式启用：

\begin{verbatim}
\SynapticDefineTheme{brand}{blue!70!black}{gray!80!black}[orange]
\SynapticSetup{theme=brand}
\end{verbatim}

参数依次为小写标识符、主色、次级色与可选强调色；已有主题不可被原地覆盖。
以下语义令牌始终可用于扩展，并始终由当前主题推导：

{\setlength{\parindent}{0pt}\noindent\begin{tabularx}{\textwidth}{@{}l l Y@{}}
  \toprule
  \textbf{令牌} & \textbf{角色} & \textbf{用途} \\
  \midrule
  \texttt{syn-primary} & 品牌色 & 标题、标题规则、强调 \\
  \texttt{syn-secondary} & 墨色 & 弱化文字、注记色调 \\
  \texttt{syn-accent} & 强调色 & 可选强调 \\
  \texttt{syn-text-dark} & 正文墨色 & 可读正文 \\
  \texttt{syn-muted-text} & 可读弱化 & 辅助标签 \\
  \texttt{syn-surface} & 卡片底色 & 盒子与表面 \\
  \texttt{syn-border} & 卡片边框 & 盒子边框 \\
  \texttt{syn-rule} & 装饰发丝线 & 分隔线 \\
  \texttt{syn-danger} & 语义危险 & 红色警示 \\
  \texttt{syn-warning} & 语义警告 & 琥珀色警示 \\
  \texttt{syn-success} & 语义成功 & 练习与待办 \\
  \bottomrule
\end{tabularx}}

讲义模式提供教学卡片环境：\texttt{synlearninggoals} 与
\texttt{synlecturesummary}，以及上文描述的课程元数据命令。笔记模式提供
\texttt{synkeyidea}、\texttt{synquestion}、\texttt{syntodo} 与
\texttt{synsummary}。

\section{书籍流程与致谢}

书籍结构由三个命令明确划分：
\begin{itemize}
  \item \DocCmd{SynapticFrontMatter} 进入前置部分；
  \item \DocCmd{SynapticMainMatter} 进入正文；
  \item \DocCmd{SynapticBackMatter} 进入后置部分。
\end{itemize}
另有目录、图表目录、附录、章副标题与题词辅助命令。notes 模式兼容文档类的
双栏布局，并提供 \texttt{synkeyidea}、\texttt{synquestion}、\texttt{syntodo} 和
\texttt{synsummary}。\DocCmd{SynapticAcknowledgments} 在书籍模式中生成章，
其他模式中生成节；带星号的形式不加入目录。

\section{常见问题}

\begin{itemize}
  \item \textbf{``The synaptic package requires LuaLaTeX''} --- 您正在使用
    pdfLaTeX 或 XeLaTeX；其他引擎会被有意拒绝。
  \item \textbf{``requires a KOMA-Script document class''} --- 请在
    \texttt{scrartcl}、\texttt{scrreprt} 或 \texttt{scrbook} 下加载宏包。
  \item \textbf{``mode=book requires a chapter-based class''} --- 请使用
    \texttt{scrbook}（或其他定义了 \texttt{chapter} 与 \texttt{frontmatter}
    的 KOMA 文档类）。
  \item \textbf{``Required font ... not found''} --- 请安装文档列出的系统字体；
    Synaptic 不会换用另一套设计。
  \item \textbf{选项在加载后冻结} --- \texttt{mode}、\texttt{lang}、
    其他排版键为结构性选项，请在 \verb|\usepackage| 中设定；
    通过 \verb|\SynapticSetup| 修改会直接报错。
  \item \textbf{未找到中文字体} --- 请安装思源宋体/黑体 CN 与霞鹜文楷。
  \item \textbf{``Document title is empty''} --- 请在
    \verb|\SynapticMakeTitle| 前设置 \verb|\SynapticTitle|。
  \item \textbf{纸张尺寸不符} --- 宏包不会覆盖文档类的纸张设定，请在文档类
    选项中指定。
\end{itemize}

\section{开发命令}

\begin{verbatim}
./scripts/verify # 测试、手册、示例与诊断扫描
l3build unpack   # 从 synaptic.dtx 提取宏包文件
l3build ctan     # 创建经过测试的 TDS/CTAN 压缩包
\end{verbatim}

\SynapticAcknowledgments*
项目感谢 LaTeX3、KOMA-Script、fontspec 与 TeX 社区。

\end{document}
